
\section{Primary Evaluation: GeoNames-Scale Gazetteer Matching}

This section presents the primary evaluation of the proposed model on GeoNames-derived data.
Unlike external benchmarks, this evaluation directly reflects the intended deployment context:
large-scale gazetteer reconciliation under extreme class imbalance, with millions of candidate
pairs and high recall requirements.

\subsection{Evaluation Setup}

The model was evaluated on a held-out validation split drawn from a GeoNames-based training corpus.
Triplets were constructed to reflect realistic matching conditions, pairing true toponym variants
with large numbers of confusable negatives. Similarity was computed using cosine similarity in the
learned phonetic embedding space.

A range of similarity thresholds $\theta \in [0.5, 0.95]$ was evaluated to characterise the
precision--recall trade-off. In addition to standard F$_1$, we report F$_5$, which weights recall
five times more heavily than precision and better reflects gazetteer linking priorities.

\subsection{Dataset Statistics}

The validation split contains approximately 260k positive and 260k negative pairs, sampled from
over two million unique toponym instances. This scale is several orders of magnitude larger than
manuscript-oriented benchmarks and exposes the model to realistic noise and aliasing.

\subsection{Results by Threshold}

Table~\ref{tab:geonames_thresholds} reports performance as a function of the similarity threshold.
Lower thresholds maximise recall, while higher thresholds approach perfect precision at the cost
of reduced coverage.

\begin{table}[h]
\centering
\begin{tabular}{lcccc}
\hline
$\theta$ & Precision & Recall & F$_1$ & F$_5$ \\
\hline
0.50 & 0.957 & 0.998 & 0.977 & 0.996 \\
0.60 & 0.984 & 0.994 & 0.989 & 0.994 \\
0.70 & 0.996 & 0.985 & 0.991 & 0.986 \\
0.80 & 1.000 & 0.958 & 0.978 & 0.959 \\
0.85 & 1.000 & 0.925 & 0.961 & 0.928 \\
0.90 & 1.000 & 0.858 & 0.923 & 0.862 \\
0.95 & 1.000 & 0.699 & 0.823 & 0.707 \\
\hline
\end{tabular}
\caption{GeoNames validation performance across similarity thresholds.}
\label{tab:geonames_thresholds}
\end{table}

The best F$_5$ score is achieved at $\theta = 0.50$, reflecting the model’s ability to maintain
extremely high recall with only moderate precision loss. This behaviour is well-suited to
candidate generation pipelines, where false positives can be filtered downstream using spatial,
topological, or authority-based constraints.

\subsection{Similarity Distribution Analysis}

Figure-free analysis of similarity distributions further illustrates the model’s separation
capacity. Positive pairs exhibit a strong concentration near cosine similarity 1.0, with a
median similarity of 0.975 and a 5th percentile of 0.815. In contrast, negative pairs are centred
near zero (median 0.005), with negligible mass above 0.8 similarity.

Only 0.4\% of negative pairs exceed a similarity of 0.7, and none exceed 0.8, indicating a clean
margin between true matches and confusable negatives. Approximately 14\% of positive pairs fall
below 0.9 similarity, reflecting genuinely difficult cases such as extreme orthographic variation,
cross-lingual forms, and historical exonyms.

\subsection{Discussion}

These results demonstrate that the proposed model achieves strong separation between matching and
non-matching toponyms at GeoNames scale. The high-recall regime is particularly robust, supporting
its use as a phonetic candidate generator in large gazetteer alignment tasks. Unlike classical
string similarity measures, the learned embedding space maintains discrimination even under heavy
aliasing and multilingual variation.

\subsection{Limitations and Extensions}

While GeoNames provides broad coverage and relatively clean identifiers, it does not fully capture
the noise and inconsistency of volunteered geographic information. An OpenStreetMap-based evaluation
would introduce noisier labels and richer alias structures, and is identified as a priority for
future work.